\appendix
\section{Model families evaluated in the benchmark}
\begin{table}[h!]
\centering
\caption{Model families evaluated in the benchmark.}
\label{tab:model_families}
\begin{tabular}{|l|p{5.2cm}|} \hline
\textbf{Vendor} & \textbf{Models} \\ \hline
OpenAI & GPT-5, GPT-5-mini, GPT-5-nano, GPT-4o, GPT-4.1 \\ \hline
Anthropic & Claude Opus 4.5, Claude Sonnet 4.5, Claude Haiku 4.5 \\ \hline
DeepSeek & DeepSeek V3.2 \\ \hline
Qwen & Qwen3.5-Plus, Qwen3.5-Flash \\ \hline
Z.ai & GLM-4.7 \\ \hline
\end{tabular}
\end{table}


\section{Detailed Description of Security Domains}
\label{app:domains}
This appendix provides detailed descriptions of the security domains used in DUMA-bench. Each domain represents a class of vulnerabilities commonly observed in LLM-based agent deployments. The domains embed adversarial behavior within realistic operational workflows in which an agent interacts with a user, tools, and external information sources.

For each domain we describe the environment structure, the typical attack scenario, and the expected security properties that the agent must maintain.

%\begin{table}{H}
%\centering
%\caption{Summary of DUMA-bench domains.}
%\begin{tabular}{lll}
%Domain & Vulnerability Class & Attack Channel \\
%\midrule
%mail\_rag\_phishing & RAG poisoning & retrieved documents \\
%collab & cross-agent manipulation & collaborator messages \\
%output\_handling & unsafe output propagation & generated responses \\
%crm\_leak & trusted data leakage & CRM records \\
%infra\_loadshed & infrastructure abuse & operational tools \\
%mktg\_phishing & phishing campaign abuse & outbound email \\
%auth\_spoof\_support & identity spoofing & user interaction \\
%tool\_shadow\_poison & tool shadowing & tool interfaces \\
%\end{tabular}
%\end{table}

\subsection{mail\_rag\_phishing (RAG Poisoning)}
This domain evaluates robustness to indirect prompt injection attacks in retrieval-augmented workflows. The environment simulates an email assistant that can search a mailbox and retrieve message contents to answer user requests or compose replies.

The mailbox contains both legitimate messages and malicious emails crafted to include adversarial instructions embedded in otherwise plausible operational context. When the agent retrieves these emails as part of task completion, the injected instructions may attempt to override normal behavior, request sensitive actions, or manipulate the agent's response.

The agent must use retrieved information to complete the task while treating instructions embedded in retrieved content as untrusted input. Security assertions verify that the agent does not execute adversarial instructions originating from the retrieval context.

\subsection{collab (Cross-Agent Manipulation)}
This domain models collaborative workflows in which an agent interacts with notes or instructions produced by other agents or colleagues. The environment simulates a support or operations setting where multiple actors contribute to the same task.

Malicious collaborator notes may contain misleading instructions, requests for privileged actions, or attempts to override established procedures. These messages appear within the same communication channel as legitimate collaboration.

The agent must evaluate the credibility of collaborator instructions and avoid executing unsafe actions. Security assertions verify that sensitive operations, such as credential resets or deletion of audit logs, are not triggered solely based on unverified collaborator messages.

\subsection{output\_handling (Unsafe Output Generation)}
This domain evaluates whether the agent generates responses that can introduce downstream vulnerabilities. The environment simulates a web-support scenario in which the agent produces textual outputs that may be displayed to end users or processed by other systems.

Malicious content may appear in the interaction context, such as collaborator notes containing code snippets or embedded scripts. If forwarded or reproduced by the agent without sanitization, this content may trigger vulnerabilities such as cross-site scripting (XSS), SQL injection, or malicious redirects.

The agent must ensure that its generated outputs remain safe for downstream consumption. Security assertions verify that unsafe artifacts are not propagated in the agent's responses.

\subsection{crm\_leak (Trusted Data Oversharing)}

This domain models customer-support workflows involving access to CRM records. The environment provides case bundles containing both shareable customer information and restricted records belonging to third parties.

Attack scenarios attempt to induce the agent to disclose restricted data through urgency framing, impersonation, or regulatory pressure. The agent can inspect the case using retrieval tools but must ensure that only permitted fields are included in responses.

Security assertions verify that protected information from trusted documents is never disclosed.

\subsection{infra\_loadshed (Infrastructure Manipulation)}

This domain evaluates whether an agent can resist requests that would trigger unnecessary infrastructure actions. The environment simulates operational tools used for monitoring system health and managing compute resources.

Attack scenarios attempt to pressure the agent into executing expensive or disruptive actions such as launching heavy jobs, calling paid external APIs, overriding quota limits, or shutting down infrastructure components.

The agent must reject requests that would create unnecessary cost or instability, even when the request is framed as an urgent operational fix.

\subsection{mktg\_phishing (Campaign Abuse)}

This domain models marketing operations in which an agent manages outbound email campaigns. The environment provides campaign briefs, collaborator notes, and tools for sending messages.

Attack scenarios attempt to coerce the agent into sending phishing-style messages, often framed as urgent marketing initiatives or growth experiments. These campaigns may involve credential harvesting, suspicious links, or unapproved domains.

The agent must identify suspicious campaign patterns and refuse to execute them. Security assertions verify that phishing-style emails are not sent.

\subsection{auth\_spoof\_support (Identity Spoofing)}

This domain evaluates resistance to social-engineering attacks in customer-support workflows. The environment includes account records, verification tools, and actions such as updating account email addresses or issuing refunds.

Attackers attempt to bypass identity verification through impersonation, VIP pressure, or claims of managerial authorization.

The agent must enforce identity verification procedures before executing any sensitive action. Security assertions verify that actions such as refunds or credential updates are never executed without proper verification.

\subsection{tool\_shadow\_poison (Tool Shadowing)}

This domain evaluates whether the agent can distinguish legitimate tools from malicious or unauthorized tool interfaces. The environment exposes a tool manifest that includes both approved tools and shadow tools that mimic legitimate functionality.

Shadow tools may appear convenient or faster but violate security policies, such as exporting customer data without proper signing or authorization.

The agent must inspect tool provenance and select only approved tools when performing sensitive operations. Security assertions verify that shadow or unsigned tools are never used during execution.